\newpage
\phantomsection
\label{prf:derivative-inverse-function-at-point}

\begin{remark*}[Return]
\hyperref[thm:inverse-function-theorem-monotone]{Return to Theorem}
\end{remark*}

\begin{theorem*}[Inverse Function Theorem (Monotone Version)]
Let $I \subseteq \mathbb{R}$ be an open interval, let $f : I \to \mathbb{R}$
be strictly monotone and continuous on $I$, set $J := f(I)$, and let
$g : J \to \mathbb{R}$ be the inverse of $f$. If $f$ is differentiable
at $c \in I$ with $f'(c) \neq 0$, and if $g$ is differentiable at
$d := f(c)$, then $g'(d) = \frac{1}{f'(c)}$.
\end{theorem*}

\begin{proof}
\textbf{Professional Standard Proof.}~\\
Since $g = f^{-1}$, the composition satisfies $g(f(x)) = x$ for all $x \in I$.
Differentiating both sides at $c$ by the Chain Rule gives
\[
  g'(f(c)) \cdot f'(c) = 1,
\]
where the right-hand side is the derivative of the identity function
$x \mapsto x$ at $c$, which equals $1$. Since $f'(c) \neq 0$ by hypothesis, 
division yields
\[
  g'(d) = g'(f(c)) = \frac{1}{f'(c)}.
\]
\end{proof}

\begin{proof}
\textbf{Detailed Learning Proof.}~\\

\textbf{Step 1.} The inverse function $g$ satisfies the identity $g(f(x)) = x$ for all $x \in I$. This fundamental relationship between $f$ and its inverse $g$ provides the starting point for differentiation.

\textbf{Step 2.} Differentiate both sides of $g(f(x)) = x$ at $c$ using the Chain Rule. The left side becomes $g'(f(c)) \cdot f'(c)$, while the right side (derivative of the identity function) equals $1$. This yields the equation $g'(f(c)) \cdot f'(c) = 1$.

\textbf{Step 3.} Since $f'(c) \neq 0$ by hypothesis, division by $f'(c)$ is valid. Solving for $g'(f(c))$ gives $g'(f(c)) = \frac{1}{f'(c)}$.

\textbf{Step 4.} Since $d = f(c)$, substitution yields $g'(d) = \frac{1}{f'(c)}$, establishing the desired formula.
\end{proof}

\begin{remark*}[Proof structure]
This is a direct proof using the Chain Rule. The strategy exploits the fundamental identity $g \circ f = \text{id}$ satisfied by inverse functions. Differentiating this identity via the Chain Rule produces the multiplicative relationship $g'(f(c)) \cdot f'(c) = 1$, and the nonzero condition $f'(c) \neq 0$ makes division valid to isolate $g'(d)$. The proof assumes the differentiability of $g$ at $d$ as a hypothesis; establishing this differentiability independently requires the full constructive argument of the Inverse Function Theorem.
\end{remark*}

\begin{remark*}[Dependencies]~\\
\begin{itemize}
  \item \hyperref[thm:chain-rule]{Chain Rule}
  \item \hyperref[prop:derivative-identity]{Derivative of the Identity}
\end{itemize}
\end{remark*}